\section{Open Questions}

The following five questions remain genuinely open after this replication and are
grounded in a close re-read of Ulyanenko et al.\ 2019 against its methodological
choices and design limitations. None is answered by the published tables.

\subsection*{Q1. Is the chronic-vs-acute slope ratio driven by concurrent repair or by acute-mode signaling saturation?}
\textbf{Basis.} The paper reports $\gamma$H2AX slopes of 0.0210 (acute) vs.\ 0.0080
(chronic) foci/mGy and attributes the $\sim$2.6$\times$ ratio to concurrent DDR
resolution during chronic exposure (Section 3.1). But the same ratio would appear if
acute foci-formation kinetics saturated within the 1--10 min acute exposure window at
30--300 mGy. The paper does not test any intermediate dose rate between 0.1 and
30 mGy/min, so it cannot distinguish these two mechanisms.

\textbf{Next steps.} Run the same MSC line at 3--5 intermediate dose rates
(e.g., 0.5, 2, 8 mGy/min) at fixed cumulative 300 mGy. A saturation mechanism predicts a
sigmoidal foci-vs-rate curve; a concurrent-repair mechanism predicts a smooth
log-linear increase. Alternatively, apply a DNA-PK/ATM inhibitor (KU-55933) during
chronic exposure --- if the chronic-mode foci count rises to match acute, active repair
is confirmed; if it does not, saturation is the more likely explanation.

\subsection*{Q2. Is pATM foci formation ATM-kinase-dependent at low chronic dose?}
\textbf{Basis.} The paper treats pATM as a DSB-signaling marker, but at 30 mGy chronic
the pATM $K$-value (0.0075 foci/mGy) is comparable to $\gamma$H2AX (0.0080) while the
co-localization is reported as basal --- implying most pATM foci at the lowest chronic
dose may not co-mark DSBs. This is not called out in the paper's discussion.

\textbf{Next steps.} Repeat the low-dose chronic arm (30, 60, 100 mGy) with (a) an ATM
kinase inhibitor (KU-55933, 10 $\mu$M, 1 h pre-treatment) and (b) N-acetylcysteine ROS
scavenger pre-treatment. An ROS-driven pATM signal will drop under NAC without changing
$\gamma$H2AX; an ATM-kinase-driven signal will drop under KU-55933 with both markers
falling together.

\subsection*{Q3. Does percentage co-localization miss the biologically meaningful signal?}
\textbf{Basis.} Paper reports 43\% (acute) to 67\% (chronic-late) co-localization but
does not decompose into (i) fraction of $\gamma$H2AX$^+$ foci also pATM$^+$ vs.\
(ii) fraction of pATM$^+$ foci also $\gamma$H2AX$^+$. These are asymmetric and can move
in opposite directions across doses. The reported number is compatible with both
improving DSB signaling and increasing background pATM noise.

\textbf{Next steps.} Reanalyze the microscopy image data (if authors can share; ~3600--4800
cells/point are described) with a per-focus co-localization pipeline (ilastik +
custom Python) reporting the two asymmetric fractions separately, plus a per-cell
scatter of $\gamma$H2AX vs.\ pATM counts. A cleaner endpoint is per-cell Pearson
correlation across the population, which is direction-invariant.

\subsection*{Q4. How much of the slope variance is donor vs.\ passage vs.\ technical noise?}
\textbf{Basis.} The methods state ``primary human bone-marrow MSCs (Biolot, Russia)''
without specifying donor count, and use passage 5--6 without a passage-controlled
sub-arm. The $R^2 = 0.888$ for chronic $\gamma$H2AX (vs.\ 0.988 for acute) suggests
substantially higher scatter that could be donor-driven, passage-drift-driven, or
genuine biological noise --- the paper cannot distinguish.

\textbf{Next steps.} Design a mini-replication with 3--5 characterized MSC donors
(e.g., RoosterBio or ATCC PCS-500-012), each irradiated at 3 passages (P3, P6, P9) at
the same acute dose window. Fit a mixed-effects model with dose as fixed and
donor+passage as random; report variance components. This costs $\sim$\$8k in cells +
one operator-month and would rigorously bound the confound.

\subsection*{Q5. Do the acute-vs-chronic DDR differences translate to iPSC-derived MSCs?}
\textbf{Basis.} The paper's findings are entirely on primary bone-marrow-derived MSCs.
iPSC-derived MSCs are increasingly used in regenerative medicine and have a rejuvenated
epigenetic profile that alters heterochromatin distribution --- a key determinant of
DSB repair kinetics. The paper does not comment on translational scope.

\textbf{Next steps.} Repeat the core acute-vs-chronic comparison (0.1 vs.\ 30 mGy/min
at 100 and 300 mGy) on a well-characterized iPSC-MSC line (e.g., ATCC ACS-7010) matched
to bone-marrow MSCs from the same donor if available. Compare slope ratios and 6-h
repair kinetics. A large deviation (e.g., iPSC-MSCs showing acute-mode slopes at
chronic-mode dose-rates) would indicate the primary-MSC DDR result does not generalize
--- with direct clinical implications for iPSC-MSC-based radiotherapy adjuncts.
